\section*{\S~10 Verschiedenes}
\label{sec:verschiedenes}
\addcontentsline{toc}{section}{\S~10 Verschiedenes}

\begin{enumerate}[label=(\arabic*)]
    \item Für Themen, wie bspw. Unteroffiziers- und Feldwebelprüfungen, die hier nicht weiter vertieft werden, können durch die Leitung, den Stab, oder den erweiterten Stab eigene Regeln geschaffen werden. Diese bedürfen erst der Zustimmung der Mitgliederversammlung, wenn mindestens ein Drittel der Mitglieder dies verlangt.
    \item Der Stab und der erweiterte Stab können besondere Gremien einberufen, sollten diese benötigt werden. Über deren Zusammensetzung entscheidet das einberufende Organ. Mitglieder, die einen Mannschaftsdienstgrad inne haben, dürfen ebenfalls Mitglied eines solchen Gremiums sein.
    \item Die Person, die den Posten des Unit-Leiters innehat, schließt einen rechtlich bindenden Vertrag über das Lizenzrecht der Handbücher für die Sonder-Kampfrollen mit dem jeweiligen Autor. Hier wird für den symbolischen Betrag von 1 € das Recht erworben, die Handbücher auch bei Austritt aus der Unit weiter zu verwenden und diese zu bearbeiten und zu vervielfältigen. Ein vergleichbarer Vertrag wird bei der Wahl eines neuen Unit-Leiters zwischen dem alten und dem neuen Unit-Leiter geschlossen. Erst nach Abschluss dieses Vertrags haben die Handbücher Gültigkeit für ihre jeweilige Ausbildung.
\end{enumerate}